% problem-000404 3 双氧物种与电子结构
\section*{3. 双氧物种与电子结构}
⼰知：若 ${ \dot { \boldsymbol { y } } } = a { \boldsymbol { x } } + b _ { \dot { \boldsymbol { z } } }$ ，则 $\begin{array} { r } { \lceil \frac { \mathrm { d } y } { \mathrm { d } t } = a \frac { \mathrm { d } x } { \mathrm { d } t } \rceil } \end{array}$ 。式中，�、�均为与�⽆关的常数。

实验A. 在⽆FeP情况下，β-胡萝⼘素-间氯过氧化苯甲酸反应体系(β-CPBA)的反应机理1如下（其中β\*CPBA为反应中间物，CBA为间氯苯甲酸）：

$
\begin{array}{l l} \text {(i)} & \beta + \text {CPBA} \xleftrightarrow {K _ {i} ^ {\prime}} \beta^ {*} \text {CPBA} \\ \text {(ii)} & \beta^ {*} \text {CPBA} \xrightarrow {k _ {2} ^ {\prime}} \text {VA} + \text {CBA} \end{array}
$

实验B. 以FeP为催化剂，β-胡萝⼘素-间氯过氧化苯甲酸-铁卟啉反应体系(β-CPBA-FeP)的反应机理2如下（其中 $\mathrm { F e O P ^ { * } C B A } ,$ 、β\* ${ \mathrm { F e O P } } ^ { * } { \mathrm { C B A } } ,$ 、β\*CPBA为反应中间物）：

$
\begin{array}{l l} \text {(iii)} & \mathrm{FeP} + \mathrm{CPBA} \xleftrightarrow {K _ {1}} \mathrm{FeOP} ^ {*} \mathrm{CBA} \\ \text {(iv)} & \mathrm{FeOP} ^ {*} \mathrm{CBA} + \beta \xleftrightarrow {K _ {2}} \beta^ {*} \mathrm{FeOP} ^ {*} \mathrm{CBA} \\ \text {(v)} & \beta^ {*} \mathrm{FeOP} ^ {*} \mathrm{CBA} \xrightarrow {k _ {1}} \mathrm{VA} + \mathrm{FeP} + \mathrm{CBA} \\ \text {(vi)} & \beta + \mathrm{CPBA} \xleftrightarrow {k _ {i}} \beta^ {*} \mathrm{CPBA} \\ \text {(vii)} & \beta^ {*} \mathrm{CPBA} \xleftrightarrow {k _ {2}} \mathrm{VA} + \mathrm{CBA} \end{array} \quad \text {快速平衡}
$

对该体系的实验结果进⾏曲线拟合，可得表3-1数据 $( k _ { \mathrm { o b s } }$ 为反应的表观速率常数）；

表 3-1

\textit{（原卷此处含表格/仪器试剂表，详见 PDF 或 MD 源文件。）}

\subsection*{3-1}
对β-CBPA体系，根据实验测得的表观速率常数 $\cdot k _ { \mathrm { o b s } } ^ { \prime } ( 2 9 3 . 2 ~ \mathrm { K } ) = 4 . 7 9 5 \times 1 0 ^ { - 4 } ~ \mathrm { s } ^ { - 1 } , ~ k _ { \mathrm { o b s } } ^ { \prime } ( 3 0 1 . 2 ~ \mathrm { K } ) = 8 . 2 8 5$ $\times 1 0 ^ { - 4 } \mathrm { { s ^ { - 1 } } _ { \mathrm { { o } } } }$ 。求反应的表观活化能 $E _ { a , \mathbf { o b s } ^ { \mathsf { c } } } ^ { \prime }$

\subsection*{3-2}
根据反应机理2推导− $\cdot \frac { \mathrm { d } [ \beta ] } { \mathrm { d } t } \overset { \_ } { \ }$ [β]间关系的速率⽅程，并给出 $\updownarrow k _ { \mathrm { o b s } }$ 的表达式。

\subsection*{3-3}
已知 β-CPBA-FeP 体系反应的表观活化能 $E _ { a , 0 \mathrm { b } s } = 4 7 . 0 2$ kJ $\mathrm { m o l ^ { - 1 } }$ 和机理2中(v)步的活化能 $E _ { a , 1 } = 4 2 . 4 3$ kJ mol−1，计算(vii)步的活化能 $E _ { a , 2 \circ }$

\subsection*{3-4}
分别根据以下条件，说明β-CPBA和β-CPBA-FeP中哪⼀个体系反应更为有利。

\subsection*{3-4}
-1 $E _ { a , 1 }$ 与 $E _ { a , 2 }$ 的结果及题中所给其他数据。

\subsection*{3-4}
-2 $E _ { a , \mathbf { o b s } } ^ { \prime }$ 与 $E _ { a , \mathrm { o b s } }$ 的结果及题中所给其他数据。

实验C. 在⾦属卟啉催化氧化反应体系中加⼊⼀些含氮⼩分⼦，会加速反成。为揭⽰反应机理，研究了FeP与咪唑类(Im)含氮⼩分⼦的配位反应热⼒学。

$
\mathrm{FeP} + 2 \mathrm{Im} \rightleftharpoons \mathrm{FePIm} _ {2}
$

表 3-2

\textit{（原卷此处含表格/仪器试剂表，详见 PDF 或 MD 源文件。）}

\subsection*{3-5}
请根据表 3-2 中 293.2 K 和 301.2K 的平衡常数计算反应 1 的 $\Delta _ { \mathrm { r } } H _ { \mathrm { m } } ^ { \ominus }$ 和 $| \Delta _ { \mathrm { { r } } } S _ { \mathrm { { m } } } ^ { \ominus }$ 。⼰知 $2 7 3 { \sim } 4 0 0 \ \mathrm { K }$ 间 $\Delta C _ { p } ^ { \ominus } \ =$ $0 _ { \circ }$

\subsection*{3-6}
⽤表3-2中所给数据和3-5的计算结果分別解释温度对反应1的影响。

\subsection*{3-7}
⼰知在⼀定温度下反应⽅向会发⽣变化，请计算反应的转向温度。
